\documentclass[UTF8,12pt,a4paper]{ctexart}

\usepackage[a4paper,top=24mm,bottom=24mm,left=25mm,right=25mm]{geometry}
\usepackage{amsmath,amssymb}
\usepackage{booktabs,array,tabularx,longtable,multirow}
\usepackage{graphicx,float,caption,subcaption}
\usepackage{siunitx}
\usepackage{xcolor}
\usepackage{fancyhdr}
\usepackage{hyperref}
\usepackage{url}

\graphicspath{{figures/}{../FourBranchRev/output/}{../Localization/output/}}
\hypersetup{
  colorlinks=true,
  linkcolor=blue!55!black,
  urlcolor=blue!60!black,
  pdftitle={分布式相控阵定位系统实测实验报告},
  pdfauthor={王越、杨天宇、杨沛峰}
}
\captionsetup{font=small,labelfont=bf}
\setlength{\parindent}{2em}
\setlength{\parskip}{0.25em}
\setlength{\emergencystretch}{2em}
\renewcommand{\arraystretch}{1.18}

\definecolor{accent}{RGB}{29,78,121}
\definecolor{lightaccent}{RGB}{233,241,248}
\pagestyle{fancy}
\fancyhf{}
\fancyhead[L]{分布式相控阵定位系统实测实验报告}
\fancyhead[R]{2026-09-26}
\fancyfoot[C]{\thepage}
\setlength{\headheight}{15pt}

\newcommand{\code}[1]{\texttt{#1}}
\newcommand{\importantbox}[1]{%
  \begin{center}
  \fcolorbox{accent}{lightaccent}{%
    \begin{minipage}{0.91\textwidth}
      #1
    \end{minipage}}
  \end{center}}

\begin{document}

\begin{titlepage}
  \centering
  \vspace*{18mm}
  {\zihao{2}\bfseries 分布式相控阵使能的物联网设备高精度定位\par}
  \vspace{10mm}
  {\zihao{1}\bfseries 实测实验报告\par}
  \vspace{8mm}
  {\Large 四分支复响应恢复与分布式相控阵 AOA 定位验证\par}
  \vfill
  \begin{tabular}{rl}
    项目编号： & 202661066 \\
    项目类型： & 校级 SRTP 创新训练项目 \\
    项目负责人： & 王越 \\
    项目成员： & 王越、杨天宇、杨沛峰 \\
    指导教师： & 范伟 \\
    实测日期： & 2026 年 8 月 29 日 \\
    报告日期： & 2026 年 9 月 26 日 \\
  \end{tabular}
  \vfill
\end{titlepage}

\pagenumbering{Roman}
\begin{abstract}
本实验使用 Anritsu MS46322B 矢量网络分析仪和 32 通道、$8\times4$ 阵元相控阵，在同一信源位置下完成四个阵列位置的 On/Off、Hadamard 和 REV 测量，以验证分布式相控阵 AOA 定位链路。三种采集方式形成四条分析分支：Branch 1 复数 DFT、Branch 2 On/Off、Branch 3 Hadamard 和 Branch 4 严格功率 REV。Branch 4 在形成 $|Z_{\mathrm{REV}}|^2$ 后完全不再使用 VNA 相位，仅由功率随受控相移的变化恢复阵元相对复响应；各阵位完成宽带方向估计后，联合实现多阵位 AOA 定位。

标准 51 频点复响应分析中，Branch 1 与 Branch 4 的阵元复响应相干系数为 0.943，对齐相位均方根误差为 $15.30^\circ$；Bartlett 四阵位定位的四条分支误差分别为 1.160 m、1.172 m、1.162 m 和 1.122 m，估计位置高度集中。统一的 21 频点全组合分析覆盖 Branch 4 的全部 3--9 次相移观测子集，共 469 种复响应方法；每种方法分别接入 Bartlett 和 MUSIC，并执行 11 种多阵位定位组合。相移观测数增加时，功率 REV 的有效率和复响应一致性稳定改善。

在这一 21 频点的平行处理框架中，Bartlett/常规波束形成与单快拍形式 MUSIC 各得到 1876 条单阵位 DOA 和 5159 条定位链。四阵位定位的全部 10 个方法分组中，MUSIC 的误差中位数均低于 Bartlett，同子集配对差值中位数改善 0.031--0.058 m，显示出这批实测数据上的一致小幅收益。两种方法使用完全相同的复响应和定位几何，使这一差别可直接归于测向谱的处理方式。

总体而言，严格功率 REV 与三个复数参考分支在复响应、方向估计和 AOA 定位结果上相互印证，表明从功率观测恢复阵元相对复响应并完成分布式相控阵 AOA 定位的算法链条已经在实物系统中跑通。MUSIC 与 Bartlett 的独立处理又给出方向一致的定位结果，且 MUSIC 在本批四阵位配对中呈稳定的小幅改进。四分支共同的约 1.1 m 绝对偏差主要体现了当前仪器、几何记录及室内传播环境叠加后的端到端系统误差；它没有掩盖 REV 响应恢复与 AOA 定位链条的一致性。

\textbf{关键词：}功率 REV；分布式相控阵；AOA 定位；MUSIC；Beamforming；实测验证
\end{abstract}

\tableofcontents
\clearpage
\pagenumbering{arabic}

\section{实验目的}

本项目希望在不为每个阵元配置独立高精度相位测量链路的条件下，利用多个受控相移状态下的合成功率恢复阵元相对复响应，并完成分布式相控阵 AOA 三维定位。各阵位的 DOA 估计是形成多阵位 AOA 定位约束的中间步骤。本次实测报告的目的包括：

\begin{enumerate}
  \item 验证严格功率 REV 能否在不使用 VNA 实测相位的条件下恢复可用于 DOA 估计的阵元相对复响应；
  \item 使用复数 DFT、On/Off 和 Hadamard 三条参考分支交叉检查功率 REV；
  \item 比较四条分支在两阵位、三阵位和四阵位条件下的 AOA 定位结果；
  \item 通过全部相移观测子集穷举，判断结论是否依赖某一个特定相移组合；
  \item 将 Bartlett/常规波束形成与单快拍形式 MUSIC 作为并行测向方法，在同一批复响应上分别完成 2/3/4 阵位定位并逐项比较；
  \item 根据实测数据如实说明当前系统的精度边界和可能误差来源。
\end{enumerate}

\section{实验系统与测量场景}

\subsection{硬件与扫频设置}

实验系统的主要参数见表~\ref{tab:hardware}。

\begin{table}[H]
  \centering
  \caption{实测硬件与采集参数}
  \label{tab:hardware}
  \begin{tabularx}{\textwidth}{>{\bfseries}p{0.25\textwidth}X}
    \toprule
    项目 & 配置 \\
    \midrule
    矢量网络分析仪 & Anritsu MS46322B \\
    阵列控制模块 & Jiexi BCS3085A-32T1-B1，1 入 32 出 APM（32 路输出） \\
    实体阵列 & $4\times8$ 微带贴片天线阵列，共 32 个阵元；本文按 $x$ 方向 8 个、$z$ 方向 4 个记为 $8\times4$ \\
    测量场地 & 实验前方案记录为 A6-409 办公室内及楼道 \\
    DUT/Probe 方案 & 4×8 贴片阵列，Horn/全向天线 \\
    阵元间距 & 行向和列向均为 3.5 cm \\
    工作频段 & 实验前系统方案为 C band 4--8 GHz；本次实际扫频为 4.5--5.5 GHz，中心频率 5.0 GHz \\
    频点数与 IFBW & 2001 点，IFBW 5 kHz \\
    阵面与朝向 & 阵面平行于 $xoz$ 平面，法向朝 $+y$；被测天线朝 $-y$ \\
    校准状态 & 未进行背景相减、逐通道校准和自由空间去嵌 \\
    \bottomrule
  \end{tabularx}
\end{table}

实验前系统框图中，VNA 一端连接 1 入 32 出 APM，APM 通过 32 路 RF chain 分别驱动贴片阵元；VNA 另一端连接待测设备或探头。方案采用一个实体相控阵，通过移动阵列至少两次形成多个观测位置：每个位置分别测 DOA，再联合定位。本次实际采集进一步扩展为同一信源下的 Location1--4 四个阵位，并额外保留 Location5 单阵位数据。

测向与定位分析设置 Bartlett/常规波束形成和单快拍形式 MUSIC 两条平行路径，四分支复响应分别送入两种测向器，再使用相同的 2/3/4 阵位几何定位器。全组合配对比较统一采用 21 个频点；另以 51 频点 Bartlett 结果展示标准四分支响应和定位的细节。两种频点口径在图表中明确区分。

阵元按四个 $z$ 高度层编号，每层包含沿 $x$ 方向排列的 8 个阵元。相对于阵列几何中心，阵元坐标为
\begin{equation}
  x=(r-4.5)\times 0.035\ \mathrm{m},\qquad
  z=(2.5-c)\times 0.035\ \mathrm{m},
\end{equation}
其中 $r=1,\ldots,8$，$c=1,\ldots,4$。

\subsection{多站几何与数据范围}

Location1--4 使用同一信源位置，可组成多站定位闭环。坐标见表~\ref{tab:geometry}，原始记录由厘米统一换算为米。

\begin{table}[H]
  \centering
  \caption{四个阵列位置与信源真值}
  \label{tab:geometry}
  \begin{tabular}{lccc}
    \toprule
    对象 & $x$ (m) & $y$ (m) & $z$ (m) \\
    \midrule
    Location1 阵列中心 & 1.32 & -0.11 & 1.69 \\
    Location2 阵列中心 & 2.42 & -0.11 & 1.69 \\
    Location3 阵列中心 & 0.00 & 0.00 & 1.69 \\
    Location4 阵列中心 & -1.21 & 0.00 & 1.69 \\
    被测信源真值 & 0.00 & 3.09 & 1.50 \\
    \bottomrule
  \end{tabular}
\end{table}

Location1--4 均包含 On/Off、Hadamard 和 REV 数据。Location5 使用另一个信源位置且只有 On/Off 数据，因此只适合作为单站观察，不进入本报告的四站闭环。结果根目录中另有一组场景元数据未闭合的 Hadamard 数据，也不纳入正式统计。

\section{四分支复响应恢复方法}

\subsection{Branch 1：复数 DFT 参考}

REV 测量包含 $0^\circ,45^\circ,\ldots,315^\circ$ 八个唯一相位以及 $360^\circ$ 闭合重复点。Branch 1 对前八个唯一相位下的复数总响应 $Z(\phi_k)$ 提取一阶 DFT 分量：
\begin{equation}
  H_{\mathrm{B1}}=\frac{1}{8}\sum_{k=1}^{8}Z(\phi_k)e^{-j\phi_k}.
\end{equation}
该分支使用 VNA 相位，与 Branch 4 共用同一套原始 REV 采集，是最直接的同源复数参考。

\subsection{Branch 2：On/Off 直接测量}

每次只开启一个目标阵元，其余阵元设置为高衰减状态，直接测得每个阵元的复数 S21，形成 $H_{\mathrm{Direct}}$。该方法直观，但会受到开关隔离度、逐状态漂移和测量链路误差影响，因此作为参考而不是无误差真值。

\subsection{Branch 3：Hadamard 正交编码}

32 个阵元同时工作并依次施加 32 组 $0^\circ/180^\circ$ 正交码字。令 Hadamard 码矩阵为 $C$、阵元响应为 $h$、码字测量为 $Y$，则
\begin{equation}
  Y=Ch,\qquad H_{\mathrm{B3}}=\frac{C^{\mathrm H}Y}{32}.
\end{equation}
该分支可获得编码复响应，但也可能受到多通道同时工作时的非线性、移相误差、互耦和漂移影响。

\subsection{Branch 4：严格功率 REV}

Branch 4 首先对 REV 原始复数总响应取模平方：
\begin{equation}
  P(\phi)=|Z_{\mathrm{REV}}(\phi)|^2.
\end{equation}
从此步骤开始，算法不再访问 VNA 相位。功率模型写为
\begin{equation}
  P(\phi)=D+2\operatorname{Re}\!\left\{C e^{j\phi}\right\},
  \qquad C=B^*H,
\end{equation}
其中 $B$ 为其余阵元形成的背景场，$H$ 为目标阵元响应。对直流项、余弦项和正弦项进行最小二乘拟合后，根据二次根关系恢复相对响应 $H/(B+H)$。该响应允许存在一个全阵列公共复比例；公共增益和公共相位零点不改变阵元间相对波前，因此不影响 DOA 估计。

\importantbox{\textbf{分支隔离原则：}Branch 4 的 Bartlett 和 MUSIC 两条测向路径均只使用由 $|Z_{\mathrm{REV}}|^2$ 恢复的 $H_{\mathrm{RevPower}}$。VNA 复相位仅用于复数参考分支和离线响应对照，不进入 Branch 4 的测向或定位输入。}

\section{比较指标、DOA 与定位方法}

\subsection{复响应一致性}

同一位置、同一频点的两组 32 阵元响应记为 $a$ 和 $b$。仅允许使用一个对全部阵元共有的复标量
\begin{equation}
  q=\frac{a^{\mathrm H}b}{a^{\mathrm H}a}
\end{equation}
对齐整体增益和相位零点。之后计算复相干系数、对齐相位 RMSE、复数 NMSE 和幅度 NRMSE。该操作不包含逐阵元调参。

\subsection{单阵位宽带 DOA}

一个阵位的一副相控阵给出入射方向（DOA）；两个及以上阵位的方向结果结合阵位坐标形成定位。对每组恢复的阵元复响应，分别应用 Bartlett/常规波束形成和 MUSIC。两条路径使用同一阵元有效掩码、导向矢量、方向网格和多阵位定位器；全组合公平比较在 4.5--5.5 GHz 内统一取 21 个频点。51 频点 Bartlett 结果另作标准四分支细节展示，不与 21 频点配对数据混算。

Bartlett 在每个频点计算归一化匹配功率，并跨频点作非相干平均：
\begin{equation}
  S_{\mathrm{B},f}(u)=\frac{|x_f^{\mathrm H}a_f(u)|^2}{\|x_f\|^2\|a_f(u)\|^2}.
\end{equation}

本次实测在每个阵位和频点恢复一个阵列复响应向量 $x_f$。MUSIC 用秩一矩阵 $R_f=x_fx_f^{\mathrm H}$ 的正交补定义噪声投影，
\begin{equation}
  P_{\mathrm N,f}=I-\frac{x_fx_f^{\mathrm H}}{x_f^{\mathrm H}x_f},\qquad
  S_{\mathrm{MUSIC},f}(u)=\frac{1}{a_f(u)^{\mathrm H}P_{\mathrm N,f}a_f(u)}.
\end{equation}
各频点伪谱非相干平均后选取最大值。单频点且导向向量模长固定时，该谱与 Bartlett 匹配功率是单调对应的；跨频点非线性聚合可能导致峰位差异。此处理不把九个相移状态或不同频点当作同频独立快拍，也不能证明多快拍 MUSIC 的分辨率优势。

\subsection{多站几何定位}

将至少两个不同阵位的 DOA 与已知阵位坐标结合，才能进行空间几何定位。两站定位取两条空间方向直线最近点的中点；三站和四站定位最小化估计点到全部方向线的垂距平方和，并用 Huber IRLS 降低异常射线影响。每个结果同时检查射线参数是否向前。

射线 RMS 仅反映多条射线的内部交汇一致性，不等于相对真实坐标的定位误差。带有共同方向偏差的射线可以在错误位置紧密交汇，因此报告始终同时给出真值误差。

\section{四分支复响应与双算法定位结果}

四分支复响应一致性由 51 频点标准分析给出。为同时展示两条测向路径，本节先列出 51 频点 Bartlett 的逐阵位细节和四阵位坐标，再列出统一 21 频点下 Bartlett 与 MUSIC 的四分支配对结果。两类频点口径各自成表，均来自相同的实测原始数据。

\subsection{复响应一致性}

表~\ref{tab:agreement}为标准 51 频点分析中，先对频点取中位数、再对四个阵位取中位数的结果。

\begin{table}[H]
  \centering
  \caption{四分支两两复响应一致性}
  \label{tab:agreement}
  \begin{tabular}{lrrrr}
    \toprule
    分支对 & 相干系数 & 相位 RMSE & 复数 NMSE & 幅度 NRMSE \\
    \midrule
    B1--B2 & 0.819 & $40.28^\circ$ & -4.91 dB & 0.384 \\
    B1--B3 & 0.971 & $14.65^\circ$ & -12.46 dB & 0.161 \\
    B1--B4 & 0.943 & $15.30^\circ$ & -9.97 dB & 0.257 \\
    B2--B3 & 0.827 & $37.40^\circ$ & -5.12 dB & 0.370 \\
    B2--B4 & 0.737 & $45.91^\circ$ & -3.41 dB & 0.477 \\
    B3--B4 & 0.913 & $22.44^\circ$ & -7.96 dB & 0.294 \\
    \bottomrule
  \end{tabular}
\end{table}

Branch 1 与 Branch 4 的同源对照最关键。Branch 4 主动丢弃 VNA 相位后仍取得 0.943 的相干系数和 $15.30^\circ$ 的相位 RMSE，说明功率 REV 已恢复 DOA 估计所需的主要阵元相对幅相结构。Branch 2 与其他分支的复响应一致性相对较弱，但后续 DOA 和定位仍与其他分支接近，提示绝对定位结果还受到分支之外的共同系统因素影响。

\begin{figure}[H]
  \centering
  \includegraphics[width=0.98\textwidth]{response_agreement_matrix.png}
  \caption{四分支相对复响应的一致性矩阵。对角线表示同一分支；非对角元素为四个阵位的中位统计。}
  \label{fig:agreement}
\end{figure}

\subsection{单阵位 DOA：Bartlett 细节与双算法一致性}

51 频点 Bartlett 测向下，四条分支在每个阵位的角误差非常接近，见表~\ref{tab:doa}。Location2 在所有分支中均为最大误差，Location4 均为最小误差。这种“误差主要随阵位变化、而不是随分支变化”的现象，是共同环境与几何误差的重要证据。

\begin{table}[H]
  \centering
  \caption{51 频点 Bartlett 测向的四分支单阵位 DOA 角误差（单位：度）}
  \label{tab:doa}
  \begin{tabular}{lrrrr}
    \toprule
    分支 & Location1 & Location2 & Location3 & Location4 \\
    \midrule
    B1 复数 DFT & 12.16 & 16.69 & 8.37 & 5.16 \\
    B2 On/Off & 11.75 & 16.93 & 9.42 & 5.44 \\
    B3 Hadamard & 12.07 & 16.73 & 8.71 & 5.31 \\
    B4 功率 REV & 11.75 & 16.25 & 8.37 & 5.16 \\
    \bottomrule
  \end{tabular}
\end{table}

在 21 频点的双算法配对中，B1、B2、B3 及 B4 的 8 相移组，MUSIC 与 Bartlett 的 DOA 方向差中位数分别为 $0.512^\circ$、$0.679^\circ$、$0.574^\circ$ 和 $0.699^\circ$。两种测向谱因而给出十分接近的入射方向；其定位差异主要来自宽带谱聚合后的细微峰位变化。

\subsection{四阵位定位：51 频点细节与 21 频点双算法配对}

真实信源坐标为 $(0,3.09,1.50)$ m。51 频点 Bartlett 四阵位结果见表~\ref{tab:fourstation} 和图~\ref{fig:localization}。

\begin{table}[H]
  \centering
  \caption{51 频点 Bartlett 测向的四分支四阵位定位结果}
  \label{tab:fourstation}
  \begin{tabular}{lrrrrr}
    \toprule
    分支 & $x$ (m) & $y$ (m) & $z$ (m) & 真值误差 (m) & 射线 RMS (m) \\
    \midrule
    B1 复数 DFT & -0.264 & 1.962 & 1.557 & 1.160 & 0.269 \\
    B2 On/Off & -0.259 & 1.949 & 1.571 & 1.172 & 0.283 \\
    B3 Hadamard & -0.259 & 1.958 & 1.554 & 1.162 & 0.277 \\
    B4 功率 REV & -0.263 & 2.001 & 1.564 & 1.122 & 0.263 \\
    \bottomrule
  \end{tabular}
\end{table}

\begin{figure}[H]
  \centering
  \includegraphics[width=\textwidth]{four_branch_localization.png}
  \caption{51 频点 Bartlett 测向下的四分支实测定位。四条分支的估计位置高度集中，且均相对真值表现出主要沿 $y$ 方向的共同偏差。}
  \label{fig:localization}
\end{figure}

四条分支的 $y$ 估计均位于约 1.95--2.00 m，而真值为 3.09 m；$x$ 和 $z$ 估计也高度接近。Branch 1 与 Branch 4 的三维估计坐标相差约 4 cm，远小于它们各自约 1.1 m 的真值误差。因此，当前绝对偏差不是 Branch 4 独有，不能解释为功率 REV 算法发生结构性错误。

51 频点 Bartlett 分析还包含每个分支的 6 个两站组合、4 个三站组合和 1 个四站组合，共 44 条定位结果。两站、三站、四站的分支间趋势一致。对应的 21 频点双算法四阵位结果见表~\ref{tab:branch-pair}；Branch 4 在此按全部 8 相移子集取中位数，故该行表示 9 个子集而非单次标准定位。

\begin{table}[H]
  \centering
  \caption{21 频点四分支双算法四阵位定位误差（单位：米）}
  \label{tab:branch-pair}
  \begin{tabular}{lrrr}
    \toprule
    响应分支 & 可配对方法数 & Bartlett中位误差 & MUSIC中位误差 \\
    \midrule
    B1 复数 DFT & 1 & 1.167 & 1.127 \\
    B2 On/Off & 1 & 1.183 & 1.125 \\
    B3 Hadamard & 1 & 1.171 & 1.135 \\
    B4 功率 REV（8 相移） & 9 & 1.166 & 1.116 \\
    \bottomrule
  \end{tabular}
\end{table}

四条分支在两种测向方法下都保持接近的定位量级；MUSIC 对每个分支给出略低的四阵位误差，为后续全相移子集比较提供统一参照。

\section{Branch 4 全相移子集双算法穷举}

\subsection{穷举空间与完整性}

从九次实测观测中任选 3--9 次，Branch 4 方法数为
\begin{equation}
  \sum_{k=3}^{9}\binom{9}{k}=466.
\end{equation}
加上 Branch 1--3 各一种正式方法，共 469 种复响应方法。每种方法运行 6 个两站、4 个三站和 1 个四站定位，因此共有
\begin{equation}
  469\times 11=5159
\end{equation}
条定位组合。Bartlett 与 MUSIC 各运行一遍，共形成两组逐项配对的 5159 条定位链。为保证公平比较，两种测向方法统一使用 21 个等间隔频点；该口径与前述 51 频点 Bartlett 细节结果分开报告。

每种测向方法各输出 1876 条单阵位 DOA 和 5159 条定位链，其中各有 5082 条获得数值定位。7 个三观测子集同时包含 $0^\circ$、$360^\circ$ 和一个其他相位，只有两个唯一名义相位，功率设计矩阵秩为 2，导致每种算法各有 77 条定位链不可恢复。数值定位的射线向前检查结果在第~\ref{sec:paired-algorithms}节对照。

\subsection{相移观测数对响应与定位的影响}

表~\ref{tab:exhaustive}中的 3--9 表示四个阵位各自使用的 REV 相移观测数，不是参与定位的阵位数。每个相移子集仍然在 Location1--4 分别恢复响应并进行四站定位。

\begin{table}[H]
  \centering
  \caption{Branch 4 全子集复响应与双算法四阵位定位汇总}
  \label{tab:exhaustive}
  \resizebox{\textwidth}{!}{%
  \begin{tabular}{rrrrrrrr}
    \toprule
    观测数 & 子集数 & 可恢复子集 & 有效率中位数 & 对 B1 相干 & 相位 RMSE & Bartlett误差 & MUSIC误差 \\
    \midrule
    3 & 84 & 77 & 83.48\% & 0.703 & $50.33^\circ$ & 1.123 m & 1.097 m \\
    4 & 126 & 126 & 90.40\% & 0.844 & $33.27^\circ$ & 1.157 m & 1.115 m \\
    5 & 126 & 126 & 93.27\% & 0.893 & $26.17^\circ$ & 1.162 m & 1.118 m \\
    6 & 84 & 84 & 94.87\% & 0.921 & $21.43^\circ$ & 1.163 m & 1.117 m \\
    7 & 36 & 36 & 96.06\% & 0.938 & $18.66^\circ$ & 1.165 m & 1.118 m \\
    8 & 9 & 9 & 96.95\% & 0.948 & $16.62^\circ$ & 1.166 m & 1.116 m \\
    9 & 1 & 1 & 97.62\% & 0.959 & $13.57^\circ$ & 1.168 m & 1.116 m \\
    \bottomrule
  \end{tabular}}
\end{table}

\begin{figure}[H]
  \centering
  \includegraphics[width=0.96\textwidth]{dual_algorithm_four_station.png}
  \caption{同一批 Branch 4 复响应在 Bartlett 与 MUSIC 下的四阵位定位误差中位数。七个相移观测数组均呈现 MUSIC 略低的趋势。}
  \label{fig:dual-algorithm}
\end{figure}

\begin{figure}[H]
  \centering
  \includegraphics[width=\textwidth]{phase_observation_summary.png}
  \caption{Branch 4 相移观测子集的响应质量与 Bartlett 定位分布。随着观测数增加，复响应有效率和一致性持续提高；四阵位定位结果维持稳定。}
  \label{fig:phasecount}
\end{figure}

\begin{figure}[H]
  \centering
  \includegraphics[width=\textwidth]{exhaustive_comparison.png}
  \caption{469 种响应方法的质量指标与 Bartlett 定位链总体分布。相移观测数增加后，复响应一致性提升且结果分布收敛；个别低观测数子集的较小误差需结合响应质量综合判断。}
  \label{fig:exhaustive}
\end{figure}

观测数由 3 增加到 9 后，有效率、相干系数和相位 RMSE 均明显改善，说明更多功率观测能够稳定提高 REV 复响应质量。两种算法的四阵位误差分别维持在约 1.12--1.17 m 和 1.10--1.12 m；MUSIC 在各观测数组均小幅低于 Bartlett，而两条曲线整体稳定。这说明复响应质量的提高与两种测向结果的一致性可同时成立，共同偏差并非某个 REV 相移子集的特有现象。

\section{MUSIC 与 Bartlett 的逐组合配对分析}
\label{sec:paired-algorithms}

469 种复响应方法分别经 Bartlett/常规波束形成和单快拍形式 MUSIC 测向，并各自进入相同的 11 种定位组合。两种算法各有 1876 条单阵位 DOA、5159 条定位链，各有 5082 条数值有效定位；7 个不可恢复的三相移子集在每种算法下各产生 77 条无效链。表~\ref{tab:music-bartlett}按同一响应子集、同一四阵位组合逐项配对；这里的 3--9 仍指每个阵位使用的相移观测数。

\begin{table}[H]
  \centering
  \caption{21 频点 MUSIC 与 Bartlett 四阵位定位配对结果}
  \label{tab:music-bartlett}
  \resizebox{\textwidth}{!}{%
  \begin{tabular}{lrrrr}
    \toprule
    响应方法 & 可恢复子集 & Bartlett中位误差 & MUSIC中位误差 & 配对差值中位数 \\
    \midrule
    B1 复数 DFT & 1 & 1.167 m & 1.127 m & -0.040 m \\
    B2 On/Off & 1 & 1.183 m & 1.125 m & -0.058 m \\
    B3 Hadamard & 1 & 1.171 m & 1.135 m & -0.036 m \\
    B4：3 相移 & 77 & 1.123 m & 1.097 m & -0.031 m \\
    B4：4 相移 & 126 & 1.157 m & 1.115 m & -0.037 m \\
    B4：5 相移 & 126 & 1.162 m & 1.118 m & -0.042 m \\
    B4：6 相移 & 84 & 1.163 m & 1.117 m & -0.045 m \\
    B4：7 相移 & 36 & 1.165 m & 1.118 m & -0.048 m \\
    B4：8 相移 & 9 & 1.166 m & 1.116 m & -0.050 m \\
    B4：9 相移 & 1 & 1.168 m & 1.116 m & -0.052 m \\
    \bottomrule
  \end{tabular}}
\end{table}

差值定义为 MUSIC 真值误差减 Bartlett 真值误差。表中 10 个四阵位方法分组的差值中位数全部为负，改善范围为 0.031--0.058 m；其中 Branch 4 的 3--9 相移各组均保持同向改善。这是本批实测数据中有价值且一致的积极信号：在相同复响应和定位几何下，MUSIC 宽带谱处理取得了略低的典型定位误差，同时与 Bartlett 一样维持约 1.1 m 的总体量级。数值有效但射线不全向前的定位，Bartlett 为 84 条、MUSIC 为 109 条；低相移数的部分两/三阵位组合还存在异常交汇，故实际比较宜优先看四阵位配对和完整有效性指标。本次每阵位、每频点只有一个复响应向量，结论对应现有单快拍形式与宽带聚合设置；在这一明确条件下，MUSIC 表现出可复现的小幅收益，值得作为本项目的有效对照算法。

\section{误差分析与讨论}

\subsection{数据能够直接支持的判断}

现有数据能够直接支持以下判断：

\begin{enumerate}
  \item 四条分支通过不同方式获得阵元复响应，但四站估计位置高度集中；
  \item 严格功率 REV 与同源复数 DFT 具有较高复响应一致性；
  \item 四条分支在每个阵位的 DOA 误差趋势相近，误差大小主要随阵位变化；
  \item 增加 Branch 4 相移观测数显著改善复响应，却没有消除共同定位偏差；
  \item MUSIC 与 Bartlett 均完成全子集定位链；MUSIC 在 10 个四阵位方法分组的误差中位数上均取得小幅改善；
  \item 因此功率 REV 算法没有表现出独立于参考分支的结构性问题。
\end{enumerate}

\subsection{可能的共同误差来源}

结合实验记录，约 1.1 m 的共同偏差可能来自以下因素的叠加：

\begin{itemize}
  \item \textbf{硬件与射频链路误差：}32 路通道的增益、相位、开关隔离度和移相误差未进行逐通道校准；线缆、连接器和控制模块也会引入公共或通道相关误差。
  \item \textbf{尺规与坐标记录误差：}阵列中心、信源位置和高度由现场尺规记录，几何中心与电磁相位中心可能不完全一致。
  \item \textbf{阵列姿态误差：}分析使用理想阵面法向和水平/竖直方向；实际偏航、俯仰和滚转的小误差会在数米传播距离上放大为明显位置偏差。
  \item \textbf{电磁干扰与室内多径：}墙面、地面、设备和人员反射会改变阵元间相对波前；不同阵位的多径结构不同，可解释 Location2 的角误差持续偏大。
  \item \textbf{未去嵌通道与背景场：}当前数据没有背景相减和自由空间去嵌，最终结果代表“仪器--阵列--环境”整体链路，而不是纯算法误差。
  \item \textbf{频率选择性与时间漂移：}宽带平均可以降低单一频点异常的影响，但不能完全消除采集期间的链路漂移和频率相关多径。
\end{itemize}

以上因素均具有工程合理性，但当前只有一次场景下的四阵位数据，没有独立的姿态复测、逐通道标定、背景测量或跨日期重复实验，因此不能从现有结果中唯一估计各误差源的贡献。报告应使用“可能来自共同系统误差”的表述，而不能把偏差确定归因于某一个因素。

\subsection{结果边界}

本实验最重要的收获是技术路径的实物可行性、四条复响应分支的相互印证，以及两种测向处理对定位结果的一致支撑。约 1.1 m 是当前实测系统的端到端绝对误差；射线 RMS 约 0.26--0.28 m，说明多阵位方向观测能够较一致地交汇。报告同时给出真值误差，以完整呈现现有系统的定位表现。

\section{结论}

\importantbox{\textbf{核心结论：}复数 DFT、On/Off、Hadamard 和严格功率 REV 四条分支给出了高度相近的分布式相控阵 AOA 定位结果。Branch 4 完全依据功率变化恢复相对复响应，仍与同源复数参考达到 0.943 的相干系数，并完成四阵位定位，实物验证了“功率观测—REV 恢复—多阵位 AOA 定位”路径。MUSIC 与 Bartlett 均得到完整的 5159 条定位链；在 10 个四阵位方法分组中，MUSIC 的误差中位数均小幅降低。四分支约 1.1 m 的共同绝对偏差反映当前整体系统条件，定位结果在不同复响应和测向方法之间保持稳定。}

具体结论如下：

\begin{enumerate}
  \item Branch 1 与 Branch 4 的标准复响应相干系数为 0.943，对齐相位 RMSE 为 $15.30^\circ$；
  \item 四条分支四站定位误差为 1.122--1.172 m，估计位置高度集中；
  \item 469 种复响应方法和 5159 条定位链的穷举结果证明结论不依赖单一相移子集或单一站数组合；
  \item 更多相移观测持续提高功率 REV 的恢复质量，并使不同相移子集的响应表现更加稳定；
  \item MUSIC 与 Bartlett 各完成 5159 条定位链的同数据对照，MUSIC 在全部 10 个四阵位方法分组中取得 0.031--0.058 m 的配对误差中位数改善；
  \item 综合判断：功率 REV 到多阵位定位的路径已获实测验证，两种测向算法形成相互参照的完整闭环；当前绝对精度约为 1.1 m，后续应用应以这一实测量级为依据。
\end{enumerate}

\clearpage
\section{复现入口与数据文件}

标准四分支闭环入口为：
\begin{verbatim}
cd('Experiment')
analysis = run_closed_loop_analysis;
\end{verbatim}

469 种方法全子集穷举入口为：
\begin{verbatim}
cd('Experiment')
results = run_exhaustive_comparison;
\end{verbatim}

统一入口依次生成四分支复响应和 Bartlett/MUSIC 双算法全组合结果；也可从同一复响应缓存单独生成双算法配对文件：
\begin{verbatim}
cd('Experiment/Localization')
comparison = run_music_bartlett_comparison;
\end{verbatim}

本报告附图由以下脚本根据冻结 CSV 结果生成：
\begin{verbatim}
cd('Experiment/report')
generate_report_figures;
\end{verbatim}

关键数据文件包括：
\begin{itemize}
  \item \path{FourBranchRev/output/response_pair_metrics.csv}：四分支响应一致性；
  \item \path{Localization/output/doa_results.csv}：标准单阵位 DOA；
  \item \path{Localization/output/localization_2_3_4_stations.csv}：标准 44 条定位结果；
  \item \path{Localization/output/exhaustive_full_pipeline.csv}：5159 条完整定位链；
  \item \path{Localization/output/exhaustive_group_summary.csv}：按分支与相移观测数汇总。
  \item \path{Localization/output/music_vs_bartlett_doa.csv}：1876 条单阵位 DOA 的逐算法配对；
  \item \path{Localization/output/music_vs_bartlett_pipeline.csv}：5159 条定位组合的逐算法配对；
  \item \path{Localization/output/MUSIC_BARTLETT_COMPARISON_REPORT.md}：对照统计与证据边界。
\end{itemize}

\end{document}
